\section{Introduction}

Large language models (LLMs) have transformed code generation, achieving remarkable performance on benchmarks like HumanEval and MBPP. Iterative repair loops---where models refine their output based on feedback---have emerged as a key technique for improving functional correctness. The dominant paradigm, exemplified by Self-Debug, relies on execution feedback: test results and stack traces guide the model toward working solutions.

Static analysis offers an alternative signal. Tools like pylint and mypy detect potential issues---undefined variables, type mismatches, unreachable code---before execution. Unlike execution traces that show \emph{where} code fails, static warnings explain \emph{why} the code is problematic. This mechanistic feedback could enable more targeted repairs.

\textbf{The gap.} Prior work has shown static analysis improves code \emph{quality} metrics. Blyth et al.\ demonstrated reductions in security issues (40\% $\rightarrow$ 13\%) and readability problems (80\% $\rightarrow$ 11\%) when integrating static feedback into LLM pipelines. However, no study has measured the impact on \emph{functional correctness}---the pass@k metrics that matter for code generation benchmarks.

\textbf{Our question.} Does integrating static analyzer feedback into iterative LLM code repair improve pass@k compared to execution-only feedback?

\textbf{This paper.} We present a methodology study that establishes the groundwork for answering this question. We develop a hypothesis validation framework with gate-based testing, implement a pylint analysis pipeline for HumanEval, and provide baseline measurements. Our existence gate---testing whether static analyzers produce actionable warnings on benchmark code---reveals a critical methodological insight: canonical solutions exhibit only 9.15\% warning rates, making them invalid proxies for LLM-generated code.

\textbf{Contributions:}
\begin{enumerate}
    \item First baseline measurement of static warning rates on HumanEval canonical solutions (9.15\%, 23 warnings across 9 categories)
    \item Validated pylint analysis pipeline for code generation research (reusable components)
    \item Gate-based hypothesis validation methodology that correctly catches methodology limitations
    \item Identification of proxy limitation: canonical solutions $\neq$ LLM output for static analysis evaluation
\end{enumerate}

We frame this as a methodology contribution. The hypothesis that static analysis improves pass@k remains untested---but we provide the validated pipeline and baseline for future work with LLM API access.
